L'azienda Zucchetti S.P.A., tramite delle web application, fornisce numerosi servizi per la gestione del personale, tra i quali sono presenti servizi per la gestione delle paghe, di selezione e formazione delle risorse umane e di monitoraggio delle presenze. L'insieme di queste web application prende il nome di suite HR (Human Resources) di Zucchetti S.P.A.. La suite HR è resa disponibile anche su dispositivi mobili, tramite l'applicazione ZConnect, che attualmente conta un ampio bacino di utenti, con circa 720.000 utenti Android e circa 120.000 utenti iOS. La versione Android di ZConnect è sviluppata in Java e Kotlin, mentre per l'ambiente iOS è stato utilizzato Swift.

Tra i numerosi servizi offerti da ZConnect, uno di particolare importanza per questo lavoro è quello di monitoraggio delle presenze. Tale servizio permette all'utente di registrare la propria presenza sul luogo di lavoro utilizzando il dispositivo mobile personale. L'utente ha quindi la possibilità di recarsi sul posto di lavoro ed una volta entrato nell'area geografica che identifica il luogo preposto per registrare la propria presenza, può effettuare la timbratura.
Il lavoratore si può quindi trovare in due situazioni: all'interno dell'area di timbratura o all'esterno. Se si trova all'interno, ZConnect permette all'utente di registrare la propria presenza sul luogo di lavoro e di visualizzare una mappa, sulla quale è rappresentata l'area di lavoro tramite la funzionalità del geofence, il quale rappresenta un perimetro virtuale attorno a un luogo specifico \cite{GeoFence}. Se invece si trova all'esterno, il lavoratore non ha la possibilità di registrare la propria presenza sul luogo di lavoro, ma può visualizzare una mappa sulla quale è rappresentata solamente la sua posizione tramite un marker.

Zucchetti S.P.A. fornisce ZConnect ad un'azienda cliente, che per motivi di riservatezza non è possibile citare, che quindi verrà indicata come Azienda X. La collaborazione con Azienda X ha fatto emergere delle limitazioni del servizio di monitoraggio delle presenze di ZConnect.
La complessa configurazione delle aree di lavoro dell'Azienda X, che spesso si estendono su ampie porzioni del territorio, ha presentato una serie di sfide significative per i dipendenti nell'utilizzo del servizio di monitoraggio delle presenze di ZConnect. In queste aree non sono presenti edifici o strutture fisiche facilmente raggiungibili, i dipendenti quindi non hanno punti di riferimento chiari in cui registrare la loro presenza.
Questo scenario ha reso complesso per i dipendenti l'accesso alle aree di lavoro designate per la timbratura. La mancanza di edifici o riferimenti facilmente identificabili ha creato un notevole disagio, poiché gli utenti si sono trovati a dover navigare in ambienti geografici sconosciuti, spesso estesi e non facilmente riconoscibili.
Inoltre, l'applicazione ZConnect non era in grado di fornire un feedback adeguato per guidare i dipendenti verso le aree di timbratura, poiché mancava un elemento fondamentale, una mappa navigabile che mostrasse chiaramente la posizione dell'utente e identificasse l'area di lavoro più vicina. Questa mancanza di un quadro visivo completo non agevolava i dipendenti nel raggiungimento delle aree di registrazione della presenza creando in loro frustrazione nell'utilizzo di questa funzionalità.

L'obiettivo di questo lavoro di tesi è quello di risolvere le limitazioni del sistema di monitoraggio delle presenze indicate dall'Azienda X. Per risolvere queste limitazioni è stata modificata la visualizzazione ed il comportamento della mappa di ZConnect.
\'E stata modificata la visualizzazione della mappa in modo da include la posizione dell'utente e l'area di timbratura più vicina, rappresentata tramite la funzionalità del geofence. Questo risultato è stato reso possibile grazie all'implementazione di una funzionalità di zoom della mappa. Quando l'utente apre la mappa, una porzione della mappa è automaticamente ingrandita e centrata sulla sua posizione attuale, che è rappresentata da un marker. Contemporaneamente, l'area di timbratura più vicina, delimitata dal geofence, è chiaramente visibile sulla mappa. Questo approccio garantisce che l'utente possa vedere la propria posizione e l'obiettivo da raggiungere immediatamente all'apertura della mappa, semplificando il processo di navigazione.
Per agevolare il raggiungimento dell'area di timbratura è stato reso il marker che rappresenta la posizione dell'utente dinamico, facendolo muovere seguendo gli spostamenti del device dell'utente, in modo tale che il lavoratore veda in tempo reale, consultando la mappa, la direzione in cui si sta spostando. Avendo questo riscontro visivo immediato per il lavoratore risulta più semplice raggiungere l'area di registrazione delle presenze anche in ambienti a lui non familiari e privi di punti di riferimento.
\'E  stata anche implementata una funzionalità di ``follow-me" che rende più agevole l'utilizzo della mappa.
La funzionalità di ``follow-me" è stata progettata per semplificare l'esperienza dell'utente durante lo spostamento. Grazie a questa funzionalità, il lavoratore non è più costretto a interagire continuamente con la mappa durante il proprio spostamento per dover vedere la propria posizione. Invece, la mappa viene resa dinamica e gestisce automaticamente l'area visualizzata a schermo.
Con la funzionalità ``follow-me" attiva, la mappa segue in tempo reale il movimento dell'utente, garantendo che la propria posizione sia sempre al centro della porzione di mappa visualizzata.

Il processo di sviluppo del miglioramento della funzionalità di monitoraggio delle presenze ha comportato un attento lavoro di analisi, progettazione, sviluppo e test, al fine di garantire un prodotto finale affidabile e soddisfacente.
Per sviluppare questa funzionalità, mi sono unito al team di sviluppo e manutenzione di ZConnect, composto da 11 membri. Questo team è responsabile sia dello sviluppo Android che iOS. All'interno del gruppo, ho svolto l'analisi, la progettazione e lo sviluppo della funzionalità, lavorando su entrambe le piattaforme Android e iOS. Ho scritto il codice in modo indipendente e successivamente è stato oggetto di revisione da parte di alcuni colleghi del team di sviluppo. Dopo la fase di sviluppo, ho collaborato con il team di testing, diverso rispetto al team di sviluppo e manutenzione di ZConnect, per condurre una serie di test funzionali sulla nuova versione di ZConnect prodotta con questo lavoro di tesi.
Superati i test tecnici la funzionalità è stata ufficialmente rilasciata nell'applicazione ZConnect, resa disponibile agli utenti sia su piattaforma Android che iOS. Attualmente la funzionalità è utilizzata da centinaia di migliaia di persone.
Questo approccio ha garantito che il prodotto fosse pienamente funzionante, testato e pronto per l'utilizzo su larga scala.